\section{Results and Analysis}
\label{sec:results}

Main results are listed in Table~\ref{tab:main_results} and
Figure~\ref{fig:accuracy_ci}.
FTA is the main method; all other evaluated variants were not retained
(see Section~\ref{sec:ablation} and Table~\ref{tab:negative_results}).

\paragraph{Final-300.}
On the unbiased 300-example evaluation (seed~71, budget~6), FTA reaches
\textbf{86.67\%} (260/300), versus L1 83.00\% (249/300), S1 82.00\%
(246/300), and TALE 78.33\% (235/300).
Paired bootstrap 95\% confidence intervals (5\,000 resamples) for FTA
minus each baseline on this split:
\begin{itemize}[nosep,leftmargin=1.5em]
  \item vs.\ L1: $+3.67$ pp, 95\% CI $[+0.33,\;+7.00]$, bootstrap
        positive-delta fraction $0.983$
  \item vs.\ S1: $+4.67$ pp, 95\% CI $[+1.00,\;+8.33]$, bootstrap
        positive-delta fraction $0.993$
  \item vs.\ TALE: $+8.33$ pp, 95\% CI $[+5.00,\;+12.00]$, bootstrap
        positive-delta fraction $1.000$
\end{itemize}
All three CI lower bounds are strictly positive on this split.

\paragraph{Aggregate-720 (primary).}
The primary claim rests on the disjoint Aggregate-720 evaluation combining
three independent sources (seeds 41, 61, 71; comprising 300, 120, and 300
examples respectively, with no example-id or question-hash overlap).
FTA reaches \textbf{80.69\%} (581/720), versus L1 77.64\% (559/720),
S1 77.08\% (555/720), and TALE 75.14\% (541/720).
Source-stratified bootstrap 95\% confidence intervals (5\,000 resamples):
\begin{itemize}[nosep,leftmargin=1.5em]
  \item vs.\ L1: $+3.06$ pp, 95\% CI $[+0.83,\;+5.42]$, bootstrap
        positive-delta fraction $0.995$
  \item vs.\ S1: $+3.61$ pp, 95\% CI $[+1.39,\;+5.69]$, bootstrap
        positive-delta fraction $0.999$
  \item vs.\ TALE: $+5.56$ pp, 95\% CI $[+3.61,\;+7.50]$, bootstrap
        positive-delta fraction $1.000$
  \item vs.\ source-local best external: $+3.06$ pp, 95\% CI $[+1.11,\;+5.00]$,
        bootstrap positive-delta fraction $0.999$
\end{itemize}
All aggregate source-stratified CI lower bounds are strictly above zero.
The margin vs.\ best external exceeds 1 pp even at the lower CI bound.

\paragraph{Pooled ensemble baseline.}
We also evaluate a simple pooled answer ensemble over the four
already-generated answers: frontier, L1, S1, and TALE
(Table~\ref{tab:pooled_ensemble}).
The primary pooled baseline uses strict majority vote over the four answers
and falls back to the frontier answer for 2--2 ties or all-different cases;
this tie-breaker is deterministic and gold-free.
On Final-300, this Pooled-4 baseline reaches 84.33\% (253/300), while FTA
reaches 86.67\% (260/300), a $+2.33$ pp difference with paired bootstrap
95\% CI $[-0.67,\,+5.67]$.
On Aggregate-720, Pooled-4 reaches 79.86\% (575/720), while FTA reaches
80.69\% (581/720), a $+0.83$ pp difference with source-stratified 95\% CI
$[-1.11,\,+2.78]$.
An external-only L1/S1/TALE majority baseline is even closer: 85.33\% on
Final-300 and 80.28\% on Aggregate-720.
Thus, FTA remains ahead by point estimate, but the pooled-ensemble comparison
is not statistically separated and is treated as a robustness check rather
than a retained superiority claim.
The external fallback order used by LDG is deterministic and fixed for
reproducibility; it is not presented as an optimized learned ranker.
A replay sensitivity check over alternative external tie orders changes the
Aggregate-720 LDG-only accuracy from 80.00\% to 80.42\%, with frontier
fallback on all-different external ties reaching 80.83\%; these variants are
reported only as sensitivity because they were not the main method.

\paragraph{FTA vs.\ External-3 majority: selective vs.\ uniform override.}
The External-3 majority baseline (L1/S1/TALE vote) reaches 85.33\% on Final-300 and
80.28\% on Aggregate-720, placing it close to FTA by point estimate.
This proximity merits interpretation, because the two approaches differ in operating
principle: External-3 applies a uniform vote over the three external baselines on every
example, while FTA retains the frontier answer on 593 of 720 Aggregate-720 examples and
overrides only when runtime failure-trace metadata indicates structural weakness (LDG) or
unanimous high-confidence external disagreement (ECO).
FTA's value is therefore not a pooled-vote point estimate but a transparent, auditable
selective-deferral rule in which each override is attributable to a named gate and a logged
trigger field from the frontier run.
The ECO component contributes only five triggers on Aggregate-720---all recovering from
incorrect frontier answers---and should be read as a low-frequency high-precision
safeguard rather than a high-coverage source of aggregate recovery.
In this evaluation, FTA and External-3 reach similar aggregate accuracy through different
mechanisms: FTA preserves the frontier signal on most examples and defers selectively,
while External-3 discards the frontier uniformly.

\paragraph{Source-by-source breakdown.}
Table~\ref{tab:source_by_source} gives the per-source accuracy for all
evaluated methods.
Estimates are consistent across sources: FTA leads by point estimate in
every individual source split, with seed~41 showing the smallest advantage
(83.33\% vs.\ 80.33\% for all three baselines) and seed~71 showing the
largest (86.67\% vs.\ 83.00\% for L1).
The seed~61 source is substantially harder in absolute terms: frontier, L1,
S1, and TALE all fall to 57.50\%, 57.50\%, 56.67\%, and 54.17\%,
respectively.
The provider/model/budget configuration is unchanged, and the mean question
length is only modestly higher than the other sources (48.5 words vs.\ 46.1
for seed~41 and 44.6 for seed~71).
The available evidence therefore points to source-sample difficulty rather
than a method-specific run configuration change.
FTA still leads the best external on seed~61 by $+1.67$ pp, and the
source-stratified aggregate confidence interval accounts for this source
variation.

\paragraph{Win/loss/tie decomposition.}
Appendix Table~\ref{tab:wlt} gives the per-example win/loss/tie counts
versus each baseline on Final-300.
FTA wins on 19 examples vs.\ L1 (losses: 8, ties: 273),
23 examples vs.\ S1 (losses: 9, ties: 268), and 27 examples vs.\ TALE
(losses: 2, ties: 271).
The extremely low loss counts vs.\ TALE confirm that the policy's recoveries
substantially outweigh its regressions on this evaluation.

\paragraph{Gate-action decomposition.}
Table~\ref{tab:gate_actions} decomposes the main method into gate actions.
Across Aggregate-720, LDG fires on 122 examples, ECO fires on 5 examples
where LDG did not fire, and no gate fires on 593 examples.
The ECO actions are rare but highly precise: before ECO fires, the frontier
answer is incorrect on all five triggered examples; after ECO fires, all
five are correct.
The same pattern holds on Final-300, where ECO fires on three examples and
contributes three wins, zero losses, and zero ties relative to the frontier
answer and to LDG-only replay.
Accordingly, ECO should be interpreted as a low-frequency high-precision
safeguard in this evaluation, not as a standalone high-coverage component
with independent statistical reliability.
Most of the aggregate recovery mass comes from LDG.

\paragraph{Sensitivity set (supporting diagnostic only).}
Table~\ref{tab:source_by_source} includes a row for a 100-example set
(seed~31, labeled ``sensitivity-only'') that was used as a development
diagnostic during early policy calibration.
Because this set was involved in parameter selection for earlier evaluated
variants, it is excluded from the primary Aggregate-720 evidence and from the
bootstrap confidence intervals reported above.
On this set, FTA achieves 82.00\%, tied with TALE and ahead of L1 (76.00\%)
and S1 (77.00\%).
When the sensitivity set is added to the aggregate ($N = 820$), the positive
point-estimate advantage of FTA over all three external baselines is preserved,
and the overall picture is unchanged.
All retained accuracy claims in this paper are based on Final-300 and
Aggregate-720 only.
